\documentclass[11pt]{article}
\usepackage[a4paper, top=18mm, bottom=18mm, left=20mm, right=20mm]{geometry}
\usepackage[T2A]{fontenc}
\usepackage[utf8]{inputenc}
\usepackage[ukrainian]{babel}
\usepackage{graphicx}
\usepackage[export]{adjustbox}
\usepackage{amsmath}
\usepackage{amssymb}
\graphicspath{{/app/Quiz/Scripts/2023/ProgAll/15BinTree/}}
\setlength{\parindent}{0pt}
\setlength{\parskip}{6pt}
\emergencystretch=3em
\pagestyle{empty}
\begin{document}
%begin
{\centering\Large\bfseries Контрольна робота\par}
\medskip
\noindent\rule{\linewidth}{0.5pt}
\smallskip
\noindent\textbf{Варіант \No\,2}\hfill \textit{ПІБ:}\,\underline{\hspace{60mm}}
\vspace{4mm}

%beginitem
1. \textbf{Що буде виведено за виконання фрагменту коду?}

%code
#include <iostream>
using namespace std;

int h(int d){
    int x = 89;
    if (d >= 3) 
        x = 3;
    else if (d == 0)
         return 4;
    else
         return 8;
    return x;
}

int main(){
    cout << h(-1);
    return 0;
}
%code

\par\vspace{5mm}
%enditem

%beginitem
2. \textbf{Що буде виведено за виконання фрагменту коду?}

%code
print('R', end=' ')
d = ('a','b','c','d','e',0,1,2,4)
print(d[-5])
%code

\par\vspace{5mm}
%enditem

%beginitem
3. \textbf{Що буде виведено за виконання фрагменту коду?}

%code
d = ('0','1','2','3','4','5','6','7','8','9')
print(d[-7:-7:2])
%code

\par\vspace{5mm}
%enditem

%beginitem
4. \textbf{Що буде виведено за виконання фрагменту коду?}

%code
try:
    n = 13
    while n>=3:
        n += -2
    print(n, end=';')
except:
    print('Z')

%code
\par\vspace{5mm}
%enditem

%beginitem
5. \textbf{Що буде виведено за виконання фрагменту коду?}

%code
cout << (!5.0>=5 || !false== 6 || 7!=6.0);
%code
\par\vspace{5mm}
%enditem

%beginitem
6. \textbf{Що буде виведено за виконання фрагменту коду?}

%code
print(4**1.0-2)
%code
\par\vspace{5mm}
%enditem

%beginitem
7. \textbf{Що буде виведено за виконання фрагменту коду?}
%code

a,b,c=9,6,7
def h(a,b=8,c=6):
    print(a,b,c,end="")

h(3,4,c=5)
print(a,b,c)

%code
\par\vspace{5mm}
%enditem

%beginitem
8. %goodpath:Scripts/2018/Mod2018_1/03-good.txt
\textbf{Відмітити все, що є коректним оператором С++:}
( 1 ) cin>>x+3;

( 2 ) (x<y && y>0)

( 3 ) return x<=2;

( 4 ) return *;
\par\vspace{5mm}
%enditem

%beginitem
9. \textbf{Що буде виведено за виконання фрагменту коду?}
try:
    for d in range(-7, -7, 2):
        if d>=6:
            break
        print(d, end=';')
        if d>=6:
            break
    else:
        print(d,end=';')
    print(d)
except:
    print('ERR')
\par\vspace{5mm}
%enditem

%beginitem
10. \textbf{Що буде виведено за виконання фрагменту коду?}

x,y,z=3,4,5
x,y,z,z=9,z,x,y
print(x,y,z)
\par\vspace{5mm}
%enditem

%end
\newpage
\end{document}

